% GENERATED FILE. Edit canonical lessons, then run npm run generate:books.
% canonical-source-hash: fnv1a64:24ebf88ec4618a60
% canonical-lessons: TE-C215-edari, TE-C215-banda, TE-C215-maidanam, TE-C215-thiyetaru, TE-C215-stediyam

\chapter{Desert, Rock, Open ground, Theatre, Stadium}
\label{ch:te-desert-rock-open-ground-theatre-stadium}
\hlchaptermodality{215}

\begin{chapteropening}
\textbf{By the end of this chapter:} \emph{I can say a desert, a rock, an open ground, a theatre and a stadium.}

Complete the last lesson of chapter 215: I can say a desert, a rock, an open ground, a theatre and a stadium.
\end{chapteropening}

\section[\texorpdfstring{\te{ఎడారి} (eḍāri)}{ఎడారి (eḍāri)}]{\te{ఎడారి} (eḍāri) — a desert}

\label{lesson:TE-C215-edari}



\begin{quote}
Before the new one: say the Telugu for a capital city, then the Telugu for a district.
\end{quote}

\subsection*{You'll want to know: \te{ఎడారి}}
\textbf{\te{ఎడారి}} — \emph{eḍāri} — \textquotedblleft{}a desert\textquotedblright{}.

\begin{grammarlens}[title={What you've built}]
One more word for this chapter.
\end{grammarlens}

\subsection*{Your turn}
\begin{itemize}
\raggedright
  \item \emph{Say it:} \emph{eḍāri}
  \item \emph{Say it:} \emph{eḍāri}, once more
  \item \emph{Say it:} say \emph{eḍāri}
  \item \emph{Recall:} say the Telugu for a capital city, then the Telugu for a district, then say \emph{eḍāri} again
\end{itemize}

\begin{tcolorbox}[breakable,colback=teal!4,colframe=teal!35!black,title={Before you move on}]
What does \te{ఎడారి} mean? (\textquotedblleft{}A desert\textquotedblright{}.) Say it once more.
\end{tcolorbox}
\section[\texorpdfstring{\te{బండ} (baṇḍa)}{బండ (baṇḍa)}]{\te{బండ} (baṇḍa) — a rock, a boulder}

\label{lesson:TE-C215-banda}



\begin{quote}
Before the new one: say the Telugu for a district, then the Telugu for a desert.
\end{quote}

\subsection*{You'll want to know: \te{బండ}}
\textbf{\te{బండ}} — \emph{baṇḍa} — \textquotedblleft{}a rock, a boulder\textquotedblright{}.

\begin{grammarlens}[title={What you've built}]
One more word for this chapter.
\end{grammarlens}

\subsection*{Your turn}
\begin{itemize}
\raggedright
  \item \emph{Say it:} \emph{baṇḍa}
  \item \emph{Say it:} \emph{baṇḍa}, once more
  \item \emph{Say it:} say \emph{baṇḍa}
  \item \emph{Recall:} say the Telugu for a district, then the Telugu for a desert, then say \emph{baṇḍa} again
\end{itemize}

\begin{tcolorbox}[breakable,colback=teal!4,colframe=teal!35!black,title={Before you move on}]
What does \te{బండ} mean? (\textquotedblleft{}A rock, a boulder\textquotedblright{}.) Say it once more.
\end{tcolorbox}
\section[\texorpdfstring{\te{మైదానం} (maidānaṁ)}{మైదానం (maidānaṁ)}]{\te{మైదానం} (maidānaṁ) — an open ground, a plain}

\label{lesson:TE-C215-maidanam}



\begin{quote}
Before the new one: say the Telugu for a desert, then the Telugu for a rock.
\end{quote}

\subsection*{You'll want to know: \te{మైదానం}}
\textbf{\te{మైదానం}} — \emph{maidānaṁ} — \textquotedblleft{}an open ground, a plain\textquotedblright{}.

\begin{grammarlens}[title={What you've built}]
One more word for this chapter.
\end{grammarlens}

\subsection*{Your turn}
\begin{itemize}
\raggedright
  \item \emph{Say it:} \emph{maidānaṁ}
  \item \emph{Say it:} \emph{maidānaṁ}, once more
  \item \emph{Say it:} say \emph{maidānaṁ}
  \item \emph{Recall:} say the Telugu for a desert, then the Telugu for a rock, then say \emph{maidānaṁ} again
\end{itemize}

\begin{tcolorbox}[breakable,colback=teal!4,colframe=teal!35!black,title={Before you move on}]
What does \te{మైదానం} mean? (\textquotedblleft{}An open ground, a plain\textquotedblright{}.) Say it once more.
\end{tcolorbox}
\section[\texorpdfstring{\te{థియేటరు} (thiyēṭaru)}{థియేటరు (thiyēṭaru)}]{\te{థియేటరు} (thiyēṭaru) — a theatre, a cinema hall}

\label{lesson:TE-C215-thiyetaru}



\begin{quote}
Before the new one: say the Telugu for a rock, then the Telugu for an open ground.
\end{quote}

\subsection*{You'll want to know: \te{థియేటరు}}
\textbf{\te{థియేటరు}} — \emph{thiyēṭaru} — \textquotedblleft{}a theatre, a cinema hall\textquotedblright{}.

\begin{grammarlens}[title={What you've built}]
One more word for this chapter.
\end{grammarlens}

\subsection*{Your turn}
\begin{itemize}
\raggedright
  \item \emph{Say it:} \emph{thiyēṭaru}
  \item \emph{Say it:} \emph{thiyēṭaru}, once more
  \item \emph{Say it:} say \emph{thiyēṭaru}
  \item \emph{Recall:} say the Telugu for a rock, then the Telugu for an open ground, then say \emph{thiyēṭaru} again
\end{itemize}

\begin{tcolorbox}[breakable,colback=teal!4,colframe=teal!35!black,title={Before you move on}]
What does \te{థియేటరు} mean? (\textquotedblleft{}A theatre, a cinema hall\textquotedblright{}.) Say it once more.
\end{tcolorbox}
\section[\texorpdfstring{\te{స్టేడియం} (sṭēḍiyaṁ)}{స్టేడియం (sṭēḍiyaṁ)}]{\te{స్టేడియం} (sṭēḍiyaṁ) — a stadium}

\label{lesson:TE-C215-stediyam}



\begin{quote}
Before the new one: say the Telugu for an open ground, then the Telugu for a theatre.
\end{quote}

\subsection*{You'll want to know: \te{స్టేడియం}}
\textbf{\te{స్టేడియం}} — \emph{sṭēḍiyaṁ} — \textquotedblleft{}a stadium\textquotedblright{}.

\begin{grammarlens}[title={What you've built}]
One more word for this chapter.
\end{grammarlens}

\subsection*{Your turn}
\begin{itemize}
\raggedright
  \item \emph{Say it:} \emph{sṭēḍiyaṁ}
  \item \emph{Say it:} \emph{sṭēḍiyaṁ}, once more
  \item \emph{Say it:} say \emph{sṭēḍiyaṁ}
  \item \emph{Recall:} say the Telugu for an open ground, then the Telugu for a theatre, then say \emph{sṭēḍiyaṁ} again
\end{itemize}

\begin{tcolorbox}[breakable,colback=teal!4,colframe=teal!35!black,title={Before you move on}]
What does \te{స్టేడియం} mean? (\textquotedblleft{}A stadium\textquotedblright{}.) Say it once more.
\end{tcolorbox}
